\documentclass[11pt,a4paper]{article}
\usepackage{amsmath,amssymb,amsfonts,amsthm,mathtools}
\usepackage[margin=1in]{geometry}
\usepackage{microtype}
\usepackage{booktabs}
\usepackage{hyperref}
\usepackage{cleveref}
\usepackage{enumitem}
\hypersetup{colorlinks=true,linkcolor=blue,citecolor=blue,urlcolor=blue}

\newtheorem{theorem}{Theorem}[section]
\newtheorem{proposition}[theorem]{Proposition}
\newtheorem{corollary}[theorem]{Corollary}
\newtheorem{lemma}[theorem]{Lemma}
\newtheorem{definition}[theorem]{Definition}
\newtheorem{problem}[theorem]{Open Problem}
\newtheorem{hypothesis}[theorem]{Hypothesis}
\theoremstyle{remark}
\newtheorem{remark}[theorem]{Remark}

\newcommand{\C}{\mathbb C}
\newcommand{\N}{\mathbb Z_{\geq0}}
\newcommand{\Znt}{Z_{\mathrm{nt}}}
\newcommand{\ord}{\operatorname{ord}}
\newcommand{\Sing}{\operatorname{Sing}}
\newcommand{\Log}{\operatorname{Log}}

\title{\textbf{Arithmetic Singularity Transfer\\for the Shifted Prime-Zeta Series}\\[3pt]
\large Möbius Visibility, RH Rigidity, and Principal-Layer Singularities}
\author{Devin J. Wright}
\date{Version 2 Draft --- September 2026}

\begin{document}
\maketitle

\begin{abstract}
Let
\[
P(s)=\sum_p p^{-s},\qquad A^+(s)=\sum_p(p-1)^{-s},\qquad \Re s>1.
\]
We derive the exact identity
\[
A^+(s)=\sum_{k=0}^{\infty}\frac{(s)_k}{k!}P(s+k).
\]
For a non-trivial zero $\rho$ of $\zeta$, rational dilations of $\rho$ form a finite fibre
\[
\mathcal C_\lambda=\{m\ge1:m\lambda\in Z_{\rm nt}\}.
\]
We attach to this fibre the multiplicity-weighted Möbius coefficient
\[
M(\lambda)=\sum_{m\in\mathcal C_\lambda}\frac{\mu(m)}m\ord_{m\lambda}\zeta.
\]
The formal transfer map is $T(\rho,m,k)=\rho/m-k$. An unconditional affine-rigidity theorem shows that equality of two such transfer points forces equality of the translation layer and a rational-dilation relation between the underlying zeros. Under RH, every non-trivial dilation fibre is a singleton, so $M(\rho/m)=\mu(m)\ord_\rho\zeta/m$.

For the principal layer $k=0$, we give a local analytic realization under explicit branch and isolation hypotheses. In particular, under RH, every squarefree principal source $\rho/m$ is an actual logarithmic singularity of $A^+$, with coefficient $\mu(m)\ord_\rho\zeta/m$. The proof uses a local Möbius representation for $P$ and the fact that every shifted term $P(s+j)$, $j\ge1$, is holomorphic near a principal source. The deeper layers $\rho/m-k$, $k\ge1$, require continuation into the left half-plane and remain open. No implication toward or away from RH is claimed.
\end{abstract}

\section{Introduction}

The prime-zeta function
\[
P(s)=\sum_p p^{-s},\qquad \Re s>1,
\]
encodes the primes through a Dirichlet series. Its classical Möbius representation links its singular source geometry to the zeros of the Riemann zeta function. A branch-sensitive treatment is required once continuation leaves the initial half-plane.

We study the shifted prime-zeta series
\[
A^+(s)=\sum_p(p-1)^{-s},\qquad \Re s>1.
\]
The elementary factorization
\[
(p-1)^{-s}=p^{-s}(1-p^{-1})^{-s}
\]
produces an exact shifted expansion. A singular source of $P$ at $\lambda$ formally enters the $k$th shifted layer at $\lambda-k$.

The paper separates three levels:
\[
\boxed{\text{formal source}\;\longrightarrow\;\text{visible source}\;\longrightarrow\;\text{actual singularity}.}
\]
The first two are arithmetic or geometric. The third requires local analytic continuation.

The principal analytic result is deliberately narrower than a global continuation theorem: under RH and explicit local hypotheses, squarefree points $\rho/m$ are realized as logarithmic singularities of $A^+$. The layers with $k\ge1$ remain continuation problems.

\section{Exact Shifted Expansion}

\begin{theorem}[Exact shifted expansion]
For $\Re s>1$,
\[
\boxed{A^+(s)=\sum_{k=0}^{\infty}\frac{(s)_k}{k!}P(s+k).}
\tag{2.1}
\]
\end{theorem}

\begin{proof}
For every prime $p$,
\[
(p-1)^{-s}=p^{-s}(1-p^{-1})^{-s}
=p^{-s}\sum_{k=0}^{\infty}\frac{(s)_k}{k!}p^{-k}.
\]
For $\Re s>1$, the resulting double series is absolutely convergent. Indeed, the inner binomial series is absolutely convergent for every prime, and the total sum is bounded by
\[
\sum_p p^{-\sigma}(1-p^{-1})^{-\sigma}
\le 2^\sigma\sum_p p^{-\sigma}<\infty,
\qquad \sigma=\Re s>1,
\]
since $p\ge2$. Fubini's theorem therefore permits interchange of the prime and $k$ sums, giving (2.1).
\end{proof}

\section{Möbius Source Geometry}

Let
\[
Z_{\rm nt}=\{\rho\in\C:\zeta(\rho)=0,\ 0<\Re\rho<1\}.
\]
For $\lambda\in\C$, define the non-trivial dilation fibre
\[
\boxed{\mathcal C_\lambda=\{m\ge1:m\lambda\in Z_{\rm nt}\}.}
\tag{3.1}
\]
When the fibre is finite, define
\[
\boxed{M(\lambda)=\sum_{m\in\mathcal C_\lambda}\frac{\mu(m)}m\ord_{m\lambda}\zeta.}
\tag{3.2}
\]

\begin{proposition}[Finite dilation fibres]
For every $\lambda\in\C$, the set $\mathcal C_\lambda$ is finite. More precisely, if $\lambda\ne0$ and $m\in\mathcal C_\lambda$, then
\[
0<m\Re\lambda<1,
\]
so, whenever $\Re\lambda>0$,
\[
m<\frac1{\Re\lambda}.
\]
If $\Re\lambda\le0$, then $\mathcal C_\lambda=\varnothing$.
\end{proposition}

\begin{proof}
If $m\lambda\in Z_{\rm nt}$, then $0<\Re(m\lambda)<1$. Since $m>0$, this gives the stated bound and finiteness.
\end{proof}

The classical local Möbius source associated with a non-trivial zero $m\lambda$ has leading logarithmic coefficient
\[
\frac{\mu(m)}m\ord_{m\lambda}\zeta.
\]
Thus the total leading coefficient at a common source $\lambda$ is the arithmetic quantity $M(\lambda)$, once the corresponding local branches and source terms have been isolated.

For a source triple $(\rho,m,k)$ define
\[
\boxed{T(\rho,m,k)=\frac{\rho}{m}-k.}
\tag{3.3}
\]
The formal source family is therefore
\[
\frac{\rho}{m}-k,\qquad \rho\in Z_{\rm nt},\ m\ge1,\ k\ge0.
\]

\section{Affine Source Rigidity}

\begin{theorem}[Affine transfer rigidity]
For $\rho_1,\rho_2\in Z_{\rm nt}$ and $m_1,m_2\ge1$, $k_1,k_2\ge0$,
\[
\frac{\rho_1}{m_1}-k_1=\frac{\rho_2}{m_2}-k_2
\]
implies
\[
\boxed{k_1=k_2,\qquad m_2\rho_1=m_1\rho_2.}
\tag{4.1}
\]
Conversely, these two relations imply equality of the transfer points.
\end{theorem}

\begin{proof}
Taking real parts gives
\[
\frac{\Re\rho_1}{m_1}-\frac{\Re\rho_2}{m_2}=k_1-k_2.
\]
Since $0<\Re\rho_i<1$, the left side lies strictly between $-1$ and $1$. The right side is an integer. Hence $k_1=k_2$. The original equality then gives
\[
\frac{\rho_1}{m_1}=\frac{\rho_2}{m_2},
\]
which is equivalent to $m_2\rho_1=m_1\rho_2$.
\end{proof}

\begin{remark}
This theorem is unconditional. It uses only the critical strip location of non-trivial zeta zeros. A collision is a coincidence of candidate source coordinates; it is not itself an assertion that $A^+$ has a singularity there.
\end{remark}

\section{RH Rigidity and Möbius Visibility}

\begin{theorem}[Singleton dilation fibres under RH]
Assume RH. For every $\rho\in Z_{\rm nt}$ and $m\ge1$,
\[
\boxed{\mathcal C_{\rho/m}=\{m\}.}
\tag{5.1}
\]
\end{theorem}

\begin{proof}
Suppose $n\rho/m=\rho'$ for some $\rho'\in Z_{\rm nt}$. Under RH,
\[
\Re\rho=\Re\rho'=\frac12.
\]
Taking real parts yields
\[
\frac{n}{m}\frac12=\frac12,
\]
so $n=m$ and then $\rho'=\rho$.
\end{proof}

\begin{corollary}[Möbius visibility under RH]
Assume RH. For every $\rho\in Z_{\rm nt}$ and $m\ge1$,
\[
\boxed{M(\rho/m)=\frac{\mu(m)}m\ord_\rho\zeta.}
\tag{5.2}
\]
Hence the leading Möbius coefficient is nonzero precisely when $m$ is squarefree.
\end{corollary}

\section{Local Prime-Zeta Realization}

The next statement isolates the analytic input. The local continuation is formulated as a hypothesis so that branch dependence and source isolation remain explicit.

\begin{hypothesis}[Local source chart]
Fix $\lambda\in Z_{\rm nt}$. Suppose there exist a simply connected slit neighborhood $U^\ast\subset U$ of $\lambda$, a branch of the local Möbius representation of $P$ on $U^\ast$, and an integer $N\ge1$ such that
\[
P(s)=\sum_{n=1}^{N}\frac{\mu(n)}n\Log_n\zeta(ns)+R_N(s)
\tag{6.1}
\]
on $U^\ast$, where each displayed logarithm is defined on its chosen slit domain and $R_N$ is holomorphic on $U$.
Assume further that $U$ contains no source point from the finitely many displayed terms other than $\lambda$.
\end{hypothesis}

\begin{lemma}[Local leading logarithm]
Assume RH and the Local source chart at $\lambda=\rho/m$, where $\rho\in Z_{\rm nt}$ has multiplicity $r$. Then, for a compatible local branch,
\[
\boxed{P(s)=r\frac{\mu(m)}m\Log(s-\lambda)+h(s),\qquad h\in\mathcal O(U).}
\tag{6.2}
\]
\end{lemma}

\begin{proof}
Choose the chart so that $U$ contains no displayed source other than $\lambda$. A non-trivial-zero contribution from the $n$th term requires $n\lambda\in Z_{\rm nt}$. Since $\lambda=\rho/m$, RH gives
\[
\Re(n\lambda)=\frac{n}{2m}.
\]
For this to equal $1/2$, one must have $n=m$. The pole of $\zeta$ produces source points $1/n$, which are real, while $\lambda$ is non-real. Thus the unique local non-trivial-zero term is the $m$th term.

Write
\[
\zeta(z)=(z-\rho)^r g(z),\qquad g(\rho)\ne0.
\]
Then locally
\[
\Log\zeta(ms)=r\Log(s-\lambda)+\text{holomorphic term}.
\]
Multiplication by $\mu(m)/m$ yields (6.2), with all remaining terms absorbed into $h$.
\end{proof}

\section{Principal-Layer Realization}

\begin{theorem}[Principal-layer singularity transfer]
Assume RH. Let $\rho\in Z_{\rm nt}$ have multiplicity $r$, let $m\ge1$ be squarefree, and set
\[
\lambda=\frac{\rho}{m}.
\]
Assume the Local source chart at $\lambda$. Then $A^+$ has a local continuation on a slit neighborhood of $\lambda$ of the form
\[
\boxed{A^+(s)=r\frac{\mu(m)}m\Log(s-\lambda)+H(s),\qquad H\in\mathcal O(U).}
\tag{7.1}
\]
Consequently,
\[
\boxed{\lambda=\frac{\rho}{m}\in\Sing(A^+).}
\tag{7.2}
\]
\end{theorem}

\begin{proof}
By the exact expansion on $\Re s>1$ and analytic continuation through the local representation, write locally
\[
A^+(s)=P(s)+\sum_{j\ge1}\frac{(s)_j}{j!}P(s+j).
\tag{7.3}
\]
The principal term $P(s)$ has the expansion (6.2).

For every $j\ge1$ and $s$ sufficiently close to $\lambda$,
\[
\Re(s+j)>1,
\]
because $\Re\lambda=1/(2m)$. Hence each $P(s+j)$ is given by its original Dirichlet series and is holomorphic near $\lambda$.

It remains to justify local uniform convergence of the shifted tail. On a sufficiently small compact neighborhood $K$ of $\lambda$, there are constants $C,a>0$ such that
\[
\left|\frac{(s)_j}{j!}\right|\le Cj^a,
\qquad s\in K,
\]
while, for all sufficiently large $j$,
\[
|P(s+j)|\le \sum_{n\ge2}n^{-j-\inf_K\Re s}\le C'2^{-j}.
\]
Thus the shifted tail is normally convergent on $K$. It defines a holomorphic function there. Combining it with the holomorphic part of $P$ proves (7.1). Since $m$ is squarefree, $\mu(m)\ne0$, so the logarithmic term has nonzero coefficient and cannot be removed by a holomorphic function.
\end{proof}

\begin{remark}
The proof uses RH for source uniqueness and the local source chart for the prime-zeta logarithm. The higher shifted terms require no continuation of $P$ because they lie in the original half-plane near a principal source.
\end{remark}

\section{Principal-Layer Theorem}

\begin{corollary}
Under RH and the Local source chart hypotheses,
\[
\boxed{
\left\{\frac{\rho}{m}:\rho\in Z_{\rm nt},\ m\ge1,\ m\ \text{squarefree}\right\}
\subseteq \Sing(A^+).
}
\tag{8.1}
\]
For a zero $\rho$ of multiplicity $r$, the corresponding logarithmic coefficient is
\[
\boxed{r\frac{\mu(m)}m.}
\tag{8.2}
\]
\end{corollary}

\section{Deeper Shifted Layers}

For $k\ge1$, the formal transfer points are
\[
\boxed{\frac{\rho}{m}-k.}
\tag{9.1}
\]
Their real parts satisfy
\[
\Re\left(\frac{\rho}{m}-k\right)<0.
\]
At such points, terms $P(s+j)$ can reach continuation-sensitive regions. The principal-layer argument therefore gives no automatic realization theorem for $k\ge1$.

\begin{problem}[Realization of deeper layers]
Determine whether the formal sources
\[
\frac{\rho}{m}-k,\qquad k\ge1,
\]
are actual singularities of $A^+$, and construct a continuation theory sufficient to establish realization or cancellation.
\end{problem}

\section{Arithmetic Cancellation}

Away from RH, several elements of a dilation fibre can occur. Their leading logarithmic coefficients can cancel arithmetically. A synthetic example is
\[
\frac{\mu(6)}6+\frac{\mu(7)}7+\frac{\mu(42)}{42}
=\frac16-\frac17+\frac1{42}=0.
\tag{10.1}
\]
This demonstrates coefficient cancellation at the arithmetic level. It does not assert that the displayed integers form a realizable zeta-zero dilation fibre.

Under RH, the singleton-fibre theorem removes this leading-coefficient cancellation mechanism for each principal source.

\section{Computational Validation}

The computations are supplementary and independent of the analytic proofs.

For $p\le100$, $K=80$, and 50-digit arithmetic, direct evaluation and the shifted expansion agree to approximately
\[
10^{-23}
\]
at the tested points
\[
s=2+0.7i,\qquad s=1.3+2i.
\]

For the first 200 non-trivial zeros, $1\le m\le30$, and $0\le k\le20$, the tested grid contains
\[
200\cdot30\cdot21=126000
\]
transfer triples. The minimum observed full complex separation is
\[
\boxed{5.76298818597532\ldots\times10^{-4}.}
\]
Among $1\le m\le30$, there are 19 squarefree values, giving
\[
200\cdot19\cdot21=79800
\]
formally visible triples under RH.

These finite computations validate the stated formulas and implementation at the tested precision. They provide no global separation theorem.

\section{Logical Status of the Main Claims}

\begin{center}
\begin{tabular}{p{0.52\textwidth}p{0.37\textwidth}}
\toprule
\textbf{Claim} & \textbf{Status}\\
\midrule
Exact shifted expansion & Proven for $\Re s>1$.\\
Finite dilation fibres & Unconditional.\\
Affine transfer rigidity & Unconditional.\\
RH singleton fibres & Conditional on RH.\\
Möbius visibility formula & Arithmetic consequence of the fibre theorem.\\
Local leading logarithm & Conditional on the stated local source chart.\\
Principal-layer realization & Conditional on RH and the local source chart.\\
Deeper-layer realization ($k\ge1$) & Open.\\
Complete local-germ cancellation & Open.\\
Any implication toward or away from RH & Open.\\
\bottomrule
\end{tabular}
\end{center}

\section{Discussion and Scope}

The paper establishes a compact chain:
\[
\boxed{
A^+
\xrightarrow{\text{exact expansion}}
P(s+k)
\xrightarrow{\text{Möbius sources}}
\frac{\rho}{m}-k
}
\]
followed by the unconditional affine rigidity relation
\[
\boxed{
\frac{\rho_1}{m_1}-k_1=\frac{\rho_2}{m_2}-k_2
\Longrightarrow
k_1=k_2,\quad m_2\rho_1=m_1\rho_2.
}
\]
Under RH,
\[
\boxed{
\mathcal C_{\rho/m}=\{m\},\qquad
M(\rho/m)=\frac{\mu(m)}m\ord_\rho\zeta.
}
\]
The principal layer then admits an actual singularity theorem after the local source chart has been fixed.

The resulting frontier is
\[
\boxed{
\underbrace{\rho/m-k}_{\text{formal source for all }k\ge0}
\qquad
\begin{cases}
\text{actual singularity under the principal-layer hypotheses},&k=0,\\
\text{continuation problem},&k\ge1.
\end{cases}}
\]

The manuscript makes no claim that the shifted series encodes an equivalent formulation of RH. A failure of RH would require an additional arithmetic bridge producing a relevant dilation relation, followed by analytic survival of the resulting source. Those bridges remain open.

\section{Conclusion}

The central structural result is the separation of source geometry from analytic realization. The affine map
\[
(\rho,m,k)\mapsto \rho/m-k
\]
has a rigid collision structure throughout the non-trivial zero sector. RH sharpens this structure to singleton dilation fibres and explicit Möbius visibility. For the principal layer, the local prime-zeta logarithm transfers directly to $A^+$, producing genuine singularities at squarefree dilates $\rho/m$ under the stated local hypotheses.

The deeper layers remain a continuation problem. Complete local-germ cancellation, global continuation, and any converse relation with RH remain outside the present theorem package.

\end{document}
